\scanpage{219}
\mspage{168}

est à la fois cartésien et cocartésien. De plus, si $W$ est un localiseur fondamental \unsure{dans|de} Cat, satisfaisant \deleted{\unsure{à}} $W(7\ \mathrm{bis})$ (en plus de $W(1,2,\unsure{3|13|{}'3})$), alors le \deleted{diagramme} carré $X/Q$ dans Cat est $W$-cocartésien.

Ce dernier énoncé résulte de ce qui précède et du cor. 2 du th. 4, compte tenu que $X/q$ est un iso. local donc est $W$-lisse.

\uline{Corollaire} : On a \inserted{(moyennant $W(7\ \mathrm{bis})$)} les équivalences

\[
\begin{aligned}
i\in W_X &\Longleftrightarrow i'\in W_X\\
p\in W_X &\Longleftrightarrow q\in W_X ,
\end{aligned}
\]

% slanted
\margin{\uline{NB} \unsure{Si on veut} seulement $i\in W_X\Rightarrow i'\in W_X$ $p\in W_X\Rightarrow q\in W_X$, il suffit que l'on ait $W(7\ \mathrm{bis})$ \uline{faible} i.e. $i_0\mathrel{\unsure{\in|\leftarrow}}W\Rightarrow i\in W$ dans l'énoncé de $W(7\ \mathrm{bis})$ \unsure{p. 76|p. 75} (Loc$(6'\mathrm a)$ avec les notations de $\underline{\overline{\mathrm{XVI}}}$)}

\[
W_X\subset\operatorname{Fl}(X^\wedge)
\]

est \unsure{défini|définie} par (cf. P. S. \ldots\ )

\[
\begin{aligned}
u\in W_X &\Longleftrightarrow X/u\colon X/F\unsure{-|\longrightarrow}X/G\\
&\hspace{4em}\text{est dans }W.
\end{aligned}
\]

Ainsi, dans une catégorie $X^\wedge$, les sommes amalgamées impliquant un mono se comportent de façon parfaite \unsure{du|au} point de vue homotopique, plus exactement pour les localiseurs $W_X$, pourvu
% archive sheet number: 219
